\section{Introdução}
\label{sec:introducao}

Os sistemas ciberfísicos (\cite{lee2017embedded}) integram a computação com processos físicos. Esses sistemas trabalham na interseção entre componentes computacionais, como os microcontroladores, e elementos que interagem com o mundo físico, como sensores e atuadores. Desde seu desenvolvimento eles têm sido empregados no desenvolvimento de diversas tecnologias, incluindo veículos autônomos, infraestruturas críticas inteligentes e, mais recentemente, sistemas médicos ciberfísicos (\textit{Medical Cyber-Physical Systems} - MCPS) (\cite{dey2018medical}). Com o avanço da biomedicina, o interesse no desenvolvimento de MCPS tem crescido, impulsionando o desenvolvimento e implementação de tecnologias de controle. Exemplos dessas aplicações incluem o pâncreas artificial (\cite{brown2019six}) e os controladores para reanimação volêmica em casos de hipovolemia (\cite{kramer2008closed}). Nesse contexto, a administração automatizada de medicamentos, na qual algoritmos de controle são empregados para ajustar continuamente a dosagem administrada ao paciente, torna-se uma das aplicações mais relevantes e promissoras dos MCPS.

O campo da farmacologia clínica é de especial interesse devido ao potencial dessas tecnologias para aumentar a eficácia terapêutica e diminuir os riscos associados à administração inadequada de medicamentos. A farmacocinética, que estuda a evolução da concentração de medicamentos no organismo em função do tempo e da dose (\cite{bailey2005drug}), permite representar esses processos através de modelos compartimentais. Com base nesses modelos, diversas estratégias de controle têm sido propostas para a administração automatizada de medicamentos (\cite{gholami2013optimal,nogueira2016positive,padmanabhan2019optimal}).

Uma alternativa para o projeto de controladores (ou, rastreadores de referência) aplicados a processos farmacocinéticos consiste no uso de técnicas de controle robusto baseadas em desigualdades matriciais lineares (\textit{Linear Matrix Inequalities} - LMIs) (\cite{morais_reference,marcomini2025projeto}). Essas técnicas permitem formular condições de estabilidade para famílias de sistemas afetados por incertezas paramétricas e variantes no tempo, tornando o rastreador válido mesmo que a modelagem física do indivíduo fosse imprecisa ou sofresse pequenas flutuações ao longo do tempo. 

A validação de estratégias de controle propostas na literatura é frequentemente realizada em ambientes de simulação. Contudo, com o objetivo de aproximar a análise do rastreador de condições práticas de operação, também têm sido desenvolvidas plataformas físicas  de emulação utilizando circuitos eletrônicos (\cite{costa2024designing}) e sistemas embarcados (\cite{desouza2025embedded, oliveira2024preliminary}).

Neste trabalho, considera-se um modelo farmacocinético discreto de três compartimentos, que representa a distribuição de um fármaco entre músculo, sangue e órgãos. Com base nesse modelo, estabelece-se uma malha fechada de controle proporcional-integral (PI) destinada ao rastreamento de uma referência de concentração de fármaco. O projeto dos ganhos do controlador é realizado através de um procedimento iterativo baseado em LMIs (\cite{morais_reference}), adaptado à malha fechada adotada. Para avaliar o comportamento do rastreador frente às incertezas paramétricas consideradas, são realizadas simulações em MATLAB/Simulink (\cite{matlab2016a}) para diferentes combinações dos parâmetros da planta. Posteriormente, a malha é implementada em dois microcontroladores ESP32, emulando respectivamente o controlador e a planta, com comunicação por UART e monitoramento por MQTT.

\newpage